\begingroup
\renewcommand{\arraystretch}{1.45}
\clearpage\fancyhead[L]{\small BESTIMMUNG DER ANZIEHUNG}\thispagestyle{empty}\pdfbookmark[0]{Bestimmung der Anziehung}{attraction-title}\SourceLeaf{340}
\vspace*{1.5em}
\begin{center}
\textbf{BESTIMMUNG DER ANZIEHUNG}

AUF EINEN BELIEBIGEN PUNKT GEGEBENER LAGE,

\textbf{DIE EIN PLANET AUSÜBEN WÜRDE, WENN SEINE MASSE}

ÜBER SEINE GESAMTE BAHN

\textbf{IM VERHÄLTNIS DER ZEIT, IN DER DIE EINZELNEN TEILE DURCHLAUFEN WERDEN,}

GLEICHMÄSSIG VERTEILT WÄRE

\bigskip
VON

\bigskip
CARL FRIEDRICH GAUSS

\bigskip
DER KÖNIGLICHEN GESELLSCHAFT DER WISSENSCHAFTEN 1818, AM 17. JANUAR, VORGELEGT.

\bigskip
\rule{0.7\linewidth}{0.4pt}

Commentationes societatis regiae scientiarum Gottingensis recentiores. Vol.~\textsc{iv}.

Göttingen \textsc{mdcccxviii}.

\rule{0.7\linewidth}{0.4pt}
\end{center}
\clearpage\fancyhead[L]{\small BESTIMMUNG DER ANZIEHUNG}\SourceLeaf{342}
{\small\begin{center}
BESTIMMUNG DER ANZIEHUNG\\
AUF EINEN BELIEBIGEN PUNKT GEGEBENER LAGE,\\
DIE EIN PLANET AUSÜBEN WÜRDE, WENN SEINE MASSE\\
ÜBER SEINE GESAMTE BAHN\\
IM VERHÄLTNIS DER ZEIT, IN DER DIE EINZELNEN TEILE DURCHLAUFEN WERDEN,\\
GLEICHMÄSSIG VERTEILT WÄRE.
\end{center}}

\begin{center}\rule{2cm}{0.3pt}\end{center}
\Needspace{23\baselineskip}
\AttractionArticle{1}
Die säkularen Veränderungen, welche die Elemente einer Planetenbahn durch die Störung eines anderen Planeten erfahren, sind von dessen Lage in seiner Bahn unabhängig und wären dieselben, ob der störende Planet sich nach \textsc{Keplers} Gesetzen in einer elliptischen Bahn bewegte oder ob man sich seine Masse so gleichmäßig über die Bahn verteilt dächte, dass den sonst in gleichen Zeitintervallen durchlaufenen Bahnteilen nun gleiche Massenteile zugewiesen würden, sofern die Umlaufzeiten des gestörten und des störenden Planeten nicht kommensurabel sind. Dieses elegante Theorem lässt sich, wenn es bisher niemand ausdrücklich ausgesprochen hat, zumindest sehr leicht aus den Grundsätzen der physischen Astronomie beweisen. Somit bietet sich ein Problem dar, das sowohl an sich als auch wegen der verschiedenen Kunstgriffe, die seine Lösung erfordert, große Aufmerksamkeit verdient: die Anziehung einer Planetenbahn oder, wenn du lieber willst. eines elliptischen Rings, dessen Dicke unendlich klein und nach dem soeben erläuterten Gesetz veränderlich ist, auf einen beliebigen Punkt gegebener Lage genau zu bestimmen.

\Needspace{40\baselineskip}
\AttractionArticle{2}
Bezeichnet man die Exzentrizität der Bahn mit $e$ und die exzentrische Anomalie eines beliebigen Punktes darin mit $E$, so wird dem Element $\mathrm{d}E$ dieser Anomalie das Element der mittleren Anomalie $(1-e\cos E)\mathrm{d}E$ entsprechen; daher wird sich das Massenelement, das dem kleinen Bahnteil zuzuweisen ist, dem jene Elemente entsprechen, zur gesamten Masse, die wir als Einheit annehmen werden, verhalten wie $(1-e\cos E)\mathrm{d}E$ zu $2\pi$, wobei $\pi$ den Halbkreisumfa\SourceLeaf{343}%
ng für den Radius 1 ausdrückt. Setzt man also den Abstand des angezogenen Punktes vom Punkt der Bahn $=\rho$, so wird die vom Bahnelement erzeugte Anziehung sein
\[
\begin{aligned}
&=\frac{(1-e\cos E)\mathrm{d}E}{2\pi\rho\rho}
\end{aligned}
\]
Wir werden die große Halbachse mit $a$ und die kleine Halbachse mit $b$ bezeichnen und jene als Abszissenlinie, den Mittelpunkt der Ellipse aber als Ursprung annehmen. Daher wird $aa-bb=aaee$ sein, die Abszisse des Bahnpunktes $=a\cos E$, die Ordinate $=b\sin E$. Schließlich werden wir den Abstand des angezogenen Punktes von der Bahnebene mit $C$ bezeichnen und die übrigen, zur großen und kleinen Achse parallelen Koordinaten mit $A$ und $B$. Nach diesen Vorbereitungen wird die Anziehung des Bahnelements in zwei zur großen und kleinen Achse parallele Anziehungen und eine dritte, zur Bahnebene senkrechte Anziehung zerlegt werden, nämlich
\[
\begin{aligned}
&\frac{(A-a\cos E)(1-e\cos E)\mathrm{d}E}{2\pi\rho^3}=\mathrm{d}\xi\\[6pt]
&\frac{(B-b\sin E)(1-e\cos E)\mathrm{d}E}{2\pi\rho^3}=\mathrm{d}\eta\\[6pt]
&\frac{C(1-e\cos E)\mathrm{d}E}{2\pi\rho^3}=\mathrm{d}\zeta
\end{aligned}
\]
wo $\rho=\surd\bigl((A-a\cos E)^2+(B-b\sin E)^2+CC\bigr)$.

Integriert man diese Differentiale von $E=0$ bis $E=360^{0}$, so werden sich die Teilanziehungen $\xi,\eta,\zeta$ in den den Koordinatenrichtungen entgegengesetzten Richtungen ergeben, aus denen sich die gesamte Anziehung zusammensetzen wird. und die sich nach der bekannten Methode auf beliebige andere Richtungen beziehen lassen.


\clearpage
\endgroup
